% P15.23: Durrett author Version 5, Section 3.2.
\chapter{Weak Convergence of Probability Laws}
\label{sec:p15-weak-convergence-of-laws}

Weak convergence describes convergence of probability laws even when
the original random variables are defined on different probability
spaces. Durrett's Section~3.2 develops a common-space quantile
construction, integral and set characterizations, and the compactness
criterion supplied by tightness
\cite[pp.~116--125]{DurrettPTEAuthor2019}.
This chapter supplies every numbered theorem and example of that
section, together with the premises needed by their proofs.
The general source presentation of the de Moivre--Laplace theorem
in Section~3.1 remains a separate treatment; elementary calculations
below establish the particular normal limits used by these examples.

\section{Distribution functions and the common-space construction}
\label{sec:p15-weak-cdf-and-coupling}

\begin{hbdefinition}{def:p15-weak-law-convergence}{Probability laws and their weak convergence}
All laws in the weak-convergence theorems are Borel probability
measures on \(\mathbb R\), and all random variables take finite
real values. For \(n\in\{1,2,\ldots\}\cup\{\infty\}\), let \(\mu_n\)
be such a law and put \(F_n(x)=\mu_n((-\infty,x])\).
Write \(\mu_n\Rightarrow\mu_\infty\), or \(F_n\Rightarrow F_\infty\),
if
\[
 F_n(x)\longrightarrow F_\infty(x)
       \quad\hbox{at every continuity point of }F_\infty.
\]
For variables \(X_n\) with these laws, the notation
\(X_n\Rightarrow X_\infty\), called \emph{convergence in
distribution}, has the same meaning. It does not require a common
probability space. Equality in distribution, written
\(X\stackrel d=Y\), means equality of laws.
A proper CDF is nondecreasing, right-continuous, takes values in
\([0,1]\), and has limits zero and one at the two infinities.
\index{weak convergence}\index{convergence in distribution}
\end{hbdefinition}

\begin{hbdefinition}{def:p15-weak-quantiles}{The source inverse convention and its plateau endpoints}
For a proper CDF \(F\) and \(u\in(0,1)\), define
\[
 Q_F(u)=\sup\{y\in\mathbb R:F(y)<u\},\qquad
 a_F(u)=Q_F(u),\qquad
 b_F(u)=\inf\{y\in\mathbb R:F(y)>u\}.
\]
The notation \(F^{-1}(u)\) in the source means \(Q_F(u)\);
no bijectivity of \(F\) is assumed. The uniform probability
space used below is
\(((0,1),\mathcal B((0,1)),\lambda)\), where \(\lambda\)
is the constructed Borel Lebesgue measure restricted to \((0,1)\).
The values at \(u=0,1\) are not part of this construction.
\end{hbdefinition}

\begin{hblemma}{thm:p15-weak-quantile-foundations}{proved=annex:p15.weak-quantile-foundations}{CDF uniqueness and measurable inverse realization}
A right-continuous nondecreasing function into \([0,1]\) has at
most countably many discontinuities. Its continuity points are
dense. A proper weak limit is unique.
For every proper CDF \(F\), \(a_F(u)\) and \(b_F(u)\) are finite,
\(a_F(u)\leq b_F(u)\), and
\[
 Q_F(u)=\inf\{y:F(y)\geq u\},\qquad
 \{u:Q_F(u)\leq x\}=\{u:u\leq F(x)\}.
\]
Consequently \(Q_F\) is Borel measurable and has CDF \(F\) on
the uniform space. In particular every proper CDF defines a
unique Borel probability law.
\end{hblemma}

\begin{hbtheorem}{thm:p15-weak-quantile-coupling}{proved=annex:p15.weak-quantile-coupling}{Durrett 3.2.8: almost-sure common-space realization}
If \(F_n\Rightarrow F_\infty\), put
\(Y_n(u)=Q_{F_n}(u)\), for \(n\in\{1,2,\ldots\}\cup\{\infty\}\),
on the uniform space. Then \(Y_n\) has CDF \(F_n\), and
\[
 Y_n(u)\longrightarrow Y_\infty(u)
       \quad\hbox{for every }u\in(0,1)\setminus E,
 \qquad
 E=\{u:a_{F_\infty}(u)<b_{F_\infty}(u)\}.
\]
The set \(E\) is at most countable and hence is a Borel null set.
Thus \(Y_n\to Y_\infty\) almost surely.
The construction preserves each individual law; it makes no
assertion that the original joint laws or independence are preserved.
\end{hbtheorem}

\section{Integral, mapping and set characterizations}
\label{sec:p15-weak-characterizations}

\begin{hbtheorem}{thm:p15-weak-bounded-continuous}{proved=annex:p15.weak-bounded-continuous}{Durrett 3.2.9: bounded continuous test functions}
For Borel probability laws on \(\mathbb R\), the following are
equivalent:
\[
 \mu_n\Rightarrow\mu_\infty;\qquad
 \int g\,d\mu_n\longrightarrow\int g\,d\mu_\infty
   \quad\hbox{for every bounded continuous }g:\mathbb R\to\mathbb R.
\]
Equivalently \(\mathbb E g(X_n)\to\mathbb E g(X_\infty)\).
These expectations are finite by boundedness.
The converse proof uses the continuous ramps
\[
 g_{x,\varepsilon}(y)=
 \begin{cases}
 1,&y\leq x,\\
 (x+\varepsilon-y)/\varepsilon,&x<y<x+\varepsilon,\\
 0,&y\geq x+\varepsilon,
 \end{cases}
 \qquad \varepsilon>0.
\]
\end{hbtheorem}

\begin{hbtheorem}{thm:p15-weak-continuous-mapping}{proved=annex:p15.weak-continuous-mapping}{Durrett 3.2.10: the exact discontinuity-set hypothesis}
Let \(g:\mathbb R\to\mathbb R\) be Borel measurable, and let
\(D_g\) be the set of points at which \(g\) is discontinuous.
The set \(D_g\) is Borel, in fact a countable union of closed sets.
If \(X_n\Rightarrow X_\infty\) and
\(\mu_\infty(D_g)=\mathbb P(X_\infty\in D_g)=0\), then
\[
 g(X_n)\Rightarrow g(X_\infty).
\]
If \(g\) is also bounded, then
\(\mathbb E g(X_n)\to\mathbb E g(X_\infty)\).
Continuity is required at almost every point for the \emph{limiting}
law. No condition that each approximating law avoid \(D_g\) is needed.
\end{hbtheorem}

\begin{hbdefinition}{def:p15-weak-boundary-sets}{Interior, closure and continuity sets}
For \(A\subseteq\mathbb R\), \(A^\circ\) is its interior,
\(\overline A\) its closure, and
\(\partial A=\overline A\setminus A^\circ\) its boundary.
A Borel set \(A\) is a \(\mu\)-continuity set when
\(\mu(\partial A)=0\).
\end{hbdefinition}

\begin{hbtheorem}{thm:p15-weak-portmanteau}{proved=annex:p15.weak-portmanteau}{Durrett 3.2.11: the portmanteau equivalences}
For Borel probability measures \(\mu_n,\mu_\infty\) on
\(\mathbb R\), the following four conditions are equivalent:
\begin{enumerate}
\item \(\mu_n\Rightarrow\mu_\infty\).
\item For every open \(G\),
\(\liminf_n\mu_n(G)\geq\mu_\infty(G)\).
\item For every closed \(K\),
\(\limsup_n\mu_n(K)\leq\mu_\infty(K)\).
\item For every Borel \(\mu_\infty\)-continuity set \(A\),
\(\mu_n(A)\to\mu_\infty(A)\).
\end{enumerate}
The open and closed sets need not be bounded. The proof uses
unit total mass when passing to complements; it is not an
assertion about arbitrary infinite measures.
\index{portmanteau theorem}
\end{hbtheorem}

\section{Helly selection, escaping mass and tightness}
\label{sec:p15-weak-selection-and-tightness}

\begin{hbdefinition}{def:p15-helly-cdf-limits}{The source's possibly defective CDF limit}
A subprobability measure is a nonnegative measure of total mass
at most one.
A \emph{Helly CDF limit} is a right-continuous nondecreasing
function \(F:\mathbb R\to[0,1]\) to which a sequence of proper
CDFs converges at every continuity point of \(F\).
Durrett denotes this convergence by \(\Rightarrow_v\) and calls
it vague convergence in this passage. That symbol here means
exactly this CDF convergence, including a possible nonzero
constant offset. The associated subprobability measure below
is obtained by subtracting the negative endpoint limit from
\(F\); the source CDF convention also records that offset.
Write
\[
 \beta=\lim_{x\to-\infty}F(x),\qquad
 \gamma=\lim_{x\to\infty}F(x).
\]
Then \(0\leq\beta\leq\gamma\leq1\). A proper limit requires
\(\beta=0,\gamma=1\).
\end{hbdefinition}

\begin{hbtheorem}{thm:p15-helly-selection}{proved=annex:p15.helly-selection}{Durrett 3.2.12: the full real-line selection construction}
Every sequence of proper CDFs has a subsequence
\(F_{n_k}\Rightarrow_v F\). The construction chooses limits
\(G(q)=\lim_kF_{n_k}(q)\) on every rational \(q\), and sets
\[
 F(x)=\inf\{G(q):q\in\mathbb Q,\ q>x\}.
\]
The resulting \(F\) has the stated range, monotonicity and right
continuity, and convergence holds at all of its continuity points.

There is a finite Borel measure \(\nu\) with total mass
\(\gamma-\beta\), uniquely specified by
\(\nu((-\infty,x])=F(x)-\beta\). For continuity points \(r<s\),
\[
 \mu_{n_k}((r,s])\longrightarrow F(s)-F(r)=\nu((r,s]).
\]
The numbers \(1-\gamma\) and \(\beta\) represent the possible
mass escaping to positive and negative infinity, respectively.
This construction creates a subprobability measure, never an
infinite measure.
\index{Helly selection theorem}\index{escaping mass}
\end{hbtheorem}

\begin{hbdefinition}{def:p15-weak-tightness}{The source tightness convention}
A sequence of proper CDFs is \emph{tight} if for every
\(\varepsilon>0\) there is a finite \(M_\varepsilon>0\) with
\[
 \limsup_{n\to\infty}
   \{1-F_n(M_\varepsilon)+F_n(-M_\varepsilon)\}
          \leq\varepsilon.
\]
The expression is the probability of
\((-\infty,-M_\varepsilon]\cup(M_\varepsilon,\infty)\);
the left endpoint is included. An equivalent uniform formulation
is that for every \(\varepsilon>0\) there exists \(R>0\) with
\(\sup_n\mu_n(\{x:|x|\geq R\})\leq\varepsilon\), as proved below.
\index{tightness}
\end{hbdefinition}

\begin{hbtheorem}{thm:p15-weak-tightness-selection}{proved=annex:p15.weak-tightness-selection}{Durrett 3.2.13: exactly when all selected limits are proper}
A sequence of proper CDFs is tight if and only if every
subsequential Helly CDF limit is proper. Equivalently, every
subsequence has a further subsequence converging weakly to a
Borel probability law. The two tightness formulations in
Definition~\ref{def:p15-weak-tightness} are equivalent.
In particular a weakly convergent sequence is tight.
\end{hbtheorem}

\begin{hbtheorem}{thm:p15-weak-coercive-tightness}{proved=annex:p15.weak-coercive-tightness}{Durrett 3.2.14: a coercive moment bound}
Let \(\varphi:\mathbb R\to[0,\infty]\) be Borel measurable and
coercive, meaning that for every \(L>0\) there is \(M>0\) such
that \(\varphi(x)\geq L\) whenever \(|x|\geq M\).
If \(C=\sup_n\int\varphi\,d\mu_n<\infty\), then the laws are tight.
For such \(L,M\), the explicit bound is
\[
 \sup_n\mu_n(\{|x|\geq M\})\leq C/L.
\]
In particular, for \(p>0\), a uniform bound on
\(\mathbb E|X_n|^p\) implies tightness.
\end{hbtheorem}

\begin{hbdefinition}{def:p15-weak-metrics}{The Levy and Ky Fan metrics}
For proper CDFs \(F,G\), define
\[
 \rho(F,G)=\inf\{\varepsilon\geq0:
 F(x-\varepsilon)-\varepsilon\leq G(x)
 \leq F(x+\varepsilon)+\varepsilon
                         \text{ for every }x\in\mathbb R\}.
\]
For finite real variables \(X,Y\) on one probability space, define
\[
 \alpha(X,Y)=\inf\{\varepsilon\geq0:
                   \mathbb P(|X-Y|>\varepsilon)\leq\varepsilon\}.
\]
For the second definition, identify variables equal almost surely.
A metric has nonnegative finite values, is symmetric, is zero
exactly on identical objects, and obeys the triangle inequality.
Convergence in a metric means that the distance to the proposed
limit tends to zero.
\index{Levy metric}\index{Ky Fan metric}
\end{hbdefinition}

\begin{hbtheorem}{thm:p15-weak-metric-premises}{proved=annex:p15.weak-metric-premises}{The source metric premises and the probability bridge}
The function \(\rho\) is a metric on proper CDFs, bounded by one,
and
\[
 \rho(F_n,F)\to0\quad\Longleftrightarrow\quad F_n\Rightarrow F.
\]
The function \(\alpha\) is a metric on finite real random
variables modulo almost-sure equality, also bounded by one.
Its defining infimum is attained, and
\[
 \alpha(X_n,X)\to0
       \quad\Longleftrightarrow\quad X_n\to X\text{ in probability}.
\]
If \(X,Y\) have CDFs \(F,G\), then
\(\rho(F,G)\leq\alpha(X,Y)\).
Consequently convergence in probability implies convergence in
distribution. Conversely, \(X_n\Rightarrow c\), for a finite
constant \(c\), implies
\(\mathbb P(|X_n-c|>\varepsilon)\to0\) for every \(\varepsilon>0\).
The metric assertions supply the full premises of the source's
argument preceding Theorem~3.2.15
\cite[Exercises~3.2.6--3.2.7 and 3.2.12, pp.~124--125]{DurrettPTEAuthor2019}.
\end{hbtheorem}

\begin{hbtheorem}{thm:p15-weak-subsequence-criterion}{proved=annex:p15.weak-subsequence-criterion}{Durrett 3.2.15 and the unique-limit method}
If every subsequence of \(X_n\) has a further subsequence
converging in distribution to \(X_\infty\), then
\(X_n\Rightarrow X_\infty\). Consequently, to prove weak
convergence it suffices to prove tightness and show that every
subsequential proper weak limit has the law of \(X_\infty\).
Both the source metric argument and a direct CDF proof are supplied.
\end{hbtheorem}

\section{Densities and the normal-limit calculations used by the examples}
\label{sec:p15-weak-density-examples}

\begin{hbdefinition}{def:p15-weak-total-variation}{The source total-variation convention}
For probability measures \(\mu,\nu\) on the same measurable
space \((S,\mathcal A)\), set
\[
 \|\mu-\nu\|_{\mathrm{TV}}
       =\sup_{A\in\mathcal A}|\mu(A)-\nu(A)|.
\]
This convention is half the integral of the absolute difference
when the two measures have densities relative to a common measure.
It is the supremum-over-events convention used in the source.
For normal examples write
\(\phi(x)=(2\pi)^{-1/2}e^{-x^2/2}\) and
\(\Phi(x)=\int_{-\infty}^x\phi(t)\,dt\).
The normalization and variance of this density have already been
proved in Theorem~\ref{thm:p15-normal-normalization}.
\end{hbdefinition}

\begin{hbtheorem}{thm:p15-scheffe-density-limit}{proved=annex:p15.scheffe-density-limit}{Scheffe's theorem in Section 3.2}
Let \((S,\mathcal A,\lambda)\) be a measure space and let
measurable \(f_n,f_\infty\geq0\) satisfy
\(\int f_n\,d\lambda=\int f_\infty\,d\lambda=1\).
If \(f_n\to f_\infty\) almost everywhere, then
\[
 \int|f_n-f_\infty|\,d\lambda
 =2\int(f_\infty-f_n)^+\,d\lambda\longrightarrow0,\qquad
 \|\mu_n-\mu_\infty\|_{\mathrm{TV}}
       =\tfrac12\int|f_n-f_\infty|\,d\lambda\longrightarrow0,
\]
where \(\mu_n(A)=\int_A f_n\,d\lambda\).
The dominating measure may have infinite total mass; all the
density measures are probabilities. On the real line this
implies weak convergence. The source assumes pointwise
convergence of Lebesgue densities, a special case of this statement.
\end{hbtheorem}

\begin{hblemma}{thm:p15-weak-elementary-asymptotics}{proved=annex:p15.weak-elementary-asymptotics}{The product limit and central-binomial premise}
Let \(\ell\in\mathbb R\) and let \(k_n\) be nonnegative integers.
If finite real arrays
\(c_{n,j}\), \(1\leq j\leq k_n\), satisfy
\[
 \max_{j\leq k_n}|c_{n,j}|\to0,\qquad
 \sum_{j=1}^{k_n}c_{n,j}\to\ell,\qquad
 \sup_n\sum_{j=1}^{k_n}|c_{n,j}|<\infty,
\]
then \(\prod_{j=1}^{k_n}(1+c_{n,j})\to e^\ell\).
Empty sums and products have values zero and one; the maximum
absolute entry of an empty row is defined to be zero.
The scalar version is also valid: if real \(c_n\to0\),
\(a_n\to\infty\), and \(a_nc_n\to\ell\in\mathbb R\), then
\((1+c_n)^{a_n}\to e^\ell\), with positive bases for all
sufficiently large \(n\).
Also, with \(C_n=4^{-n}\binom{2n}{n}\),
\[
 C_n\sqrt{\pi n}\longrightarrow1.
\]
The product argument supplies source Lemma~3.1.1 and its
Exercise~3.1.1 application. The second limit is proved directly
from cosine integrals and the established Gaussian integral,
so the examples need no assumed Stirling formula or CLT.
\end{hblemma}

\begin{hbtheorem}{thm:p15-weak-rademacher-example}{proved=annex:p15.weak-rademacher-example}{Durrett Example 3.2.1, with its needed limit proved}
Let \(\varepsilon_i\) be independent with probabilities \(1/2\)
at \(1\) and \(-1\), and put \(S_N=\sum_{i=1}^N\varepsilon_i\).
Then \(S_N/\sqrt N\Rightarrow Z\), where \(Z\) has density
\(\phi\). In particular,
\[
 \mathbb P(S_N/\sqrt N\leq x)\longrightarrow\Phi(x)
                                \qquad(x\in\mathbb R).
\]
The annex proves the uniform local binomial estimate on bounded
scaled intervals, the Riemann-sum step and the passage to odd \(N\).
This establishes the example independently of the later general
CLT presentation.
\end{hbtheorem}

\begin{hbtheorem}{thm:p15-weak-first-examples}{proved=annex:p15.weak-first-examples}{Durrett Examples 3.2.2--3.2.5}
\begin{enumerate}
\item For i.i.d. observations of any real law \(F\), the empirical
CDF \(F_n(x)=n^{-1}\sum_{i=1}^n\mathbf1_{\{X_i\leq x\}}\)
converges weakly to \(F\) on one event of probability one.
In fact the previously proved Glivenko--Cantelli theorem gives
uniform convergence, even when \(F\) has atoms.
\item If \(X\) has CDF \(F\), then \(X+1/n\Rightarrow X\).
Its CDF at \(x\) is \(F(x-1/n)\), whose limit is \(F(x-)\).
At every discontinuity of \(F\), this limit differs from \(F(x)\).
\item For \(0<p<1\), let \(X_p\) be the number of independent
Bernoulli trials, with success probability \(p\), until the first
success, assigning \(X_p=1\) on the null event of no success.
Then
\[
 \mathbb P(X_p\geq m)=(1-p)^{m-1}\ (m\geq1),\qquad
 \mathbb P(pX_p>x)=(1-p)^{\lfloor x/p\rfloor}\longrightarrow e^{-x}
 \quad(x\geq0).
\]
Thus \(pX_p\) converges in distribution, as \(p\downarrow0\),
to the exponential law with rate one.
\item For \(N\geq1\), let the observations be independent uniform
on \(\{1,\ldots,N\}\), and let \(T_N\) be the first index at
which an observation repeats an earlier one. Then
\[
 \mathbb P(T_N>m)=\prod_{j=1}^{m-1}(1-j/N)
                         \quad(1\leq m\leq N+1).
\]
For \(m=0\) the probability is one; for \(m>N+1\) it is zero.
For every \(x\geq0\),
\[
 \mathbb P(T_N/\sqrt N>x)\longrightarrow e^{-x^2/2}.
\]
The limiting CDF is zero for \(x<0\) and \(1-e^{-x^2/2}\)
for \(x\geq0\). At \(N=365,m=22\), the exact product is
approximately \(0.524305\); the source's leading limit
approximation \(e^{-22^2/(2\cdot365)}\) is approximately \(0.515296\).
\end{enumerate}
\end{hbtheorem}

\begin{hbtheorem}{thm:p15-weak-central-order-statistic}{proved=annex:p15.weak-central-order-statistic}{Durrett Example 3.2.6 and Lemma 3.2.7}
For \(n\geq1\), let \(U_1,\ldots,U_{2n+1}\) be independent
uniform variables on \((0,1)\), and let \(V_{n+1}\) be their
median, the \((n+1)\)-st smallest and also the \((n+1)\)-st
largest observation. This is a measurable random variable,
with density
\[
 f_{V_{n+1}}(x)
  =(2n+1)\binom{2n}{n}x^n(1-x)^n\mathbf1_{\{0<x<1\}}.
\]
Set \(Y_n=2\sqrt{2n}(V_{n+1}-1/2)\).
For \(|y|<\sqrt{2n}\), its density is
\[
 f_{Y_n}(y)
 =C_n(1-y^2/(2n))^n\frac{2n+1}{2n}\sqrt{\frac n2},
\]
and the density is zero outside that interval.
For every fixed real \(y\), \(f_{Y_n}(y)\to\phi(y)\).
Consequently the laws converge to the standard normal law
in total variation and hence weakly.
Density notation here denotes a density, not a singleton probability.
\end{hbtheorem}

\section{Failure boundaries and the exact continuation}
\label{sec:p15-weak-failures}

\begin{hbtheorem}{thm:p15-weak-convergence-boundaries}{proved=annex:p15.weak-convergence-boundaries}{What weak convergence does not preserve}
\begin{enumerate}
\item Point masses \(\delta_{1/n}\) converge weakly to \(\delta_0\),
but their total-variation distances from \(\delta_0\) are one.
For \(G=(0,\infty)\) and \(K=(-\infty,0]\), the portmanteau
inequalities can be strict.
\item The Borel map \(g=\mathbf1_{\{0\}}\) does not preserve this
limit: \(g(1/n)=0\), while \(g(0)=1\). Here
\(\delta_0(D_g)=1\).
\item If \(X\) is a fair sign and \(X_n=-X\) for every \(n\),
then \(X_n\Rightarrow X\), but
\(\mathbb P(|X_n-X|>1)=1\). Weak convergence on a common
space need not imply convergence in probability there.
\item A variable equal to \(n\) with probability \(1/n\) and
zero otherwise converges weakly to zero, although its expectation
is one for every \(n\). Weak convergence alone does not justify
an unbounded continuous test function.
\item The probability densities
\(f_n=\mathbf1_{[n,n+1]}\) tend pointwise to zero, whose integral
is zero. The normalization hypothesis on the limiting density
in Scheffe's theorem is necessary.
\item Let \(a,b,c\geq0\), \(a+b+c=1\), and let \(G\) be a proper
CDF. The proper CDFs
\[
 F_n(x)=a\mathbf1_{\{x\geq n\}}
       +b\mathbf1_{\{x\geq-n\}}+cG(x)
\]
converge pointwise to \(F(x)=b+cG(x)\).
Here \(\beta=b\), \(\gamma=b+c=1-a\), and the finite
remaining measure has mass \(c\). The sequence is tight
exactly when \(a=b=0\).
\end{enumerate}
\end{hbtheorem}

\begin{proofboundary}{def:p15-weak-convergence-continuation}{Source scope and the next required presentations}
The completed scope is Section~3.2's definitions, Examples
3.2.1--3.2.6, Lemma~3.2.7, Theorems~3.2.8--3.2.15 and the
unnumbered Scheffe theorem, with all their needed premises.
The distinct Section~3.1 factorial route is supplied in
\cref{thm:p15-stirling-with-remainder,thm:p15-stirling-de-moivre-laplace},
including all four source exercises. The characteristic-function
inversion, continuity and moment results through Theorem~3.3.21 are
supplied in \cref{thm:p15-characteristic-inversion-uniqueness,
thm:p15-characteristic-continuity-theorem,thm:p15-characteristic-second-moment}.
The full remaining Section~3.3 criteria, moment-problem constructions
and all section exercises are supplied in
\cref{sec:p15-characteristic-function-criteria-and-moments}.
That chapter also supplies Exercises~3.2.4, 3.2.5 and 3.2.10 through
\cref{thm:p15-criteria-weak-integration}.
The i.i.d. and general finite-row central limit theorems,
both array proof routes and all Section~3.4 exercises are supplied in
\cref{sec:p15-central-limit-theorems}, together with Erd\H{o}s--Kac
and Berry--Esseen. The complete normal-domain characterization,
including its converse and the construction of normalizing constants,
is \cref{thm:p15-nda-characterization,thm:p15-nda-normalizer-construction}.
The original truncated-mean and cutoff presentation is supplied in
\cref{thm:p15-nda-gk-truncated-mean,thm:p15-nda-gk-cutoff-route}.
The other Section~3.2 exercise presentations are supplied in
\cref{sec:p15-weak-convergence-exercises}. Together the stated
routes supply all seventeen exercises, including the Gaussian
maximum normalization, the geometric surface-measure construction
and all moment-concentration parameter cases.
\end{proofboundary}
